\documentclass[10pt]{article}

% ---------- Packages ----------

\usepackage{parskip}
\usepackage{microtype}
\usepackage{xcolor}
\usepackage{geometry}
\geometry{margin=4em, top=8em,  headheight=\paperheight}
\usepackage[export]{adjustbox}
\usepackage{array}
\usepackage{amsmath}
\usepackage{amsfonts}
\usepackage{fancyhdr}
\usepackage{luacode}
\newcommand{\assessmentType}{Respuesta Escrita}
\newcommand{\topic}{}
\pagestyle{fancy}
\fancyhf{}
\lhead{Álgebra II, \assessmentType\hspace{2em} Nombre \rule{7em}{.4pt} Apellido \rule{7em}{.4pt}
}
\chead{}
\rhead{Período \rule{3em}{.4pt}}
\usepackage{lastpage}
\usepackage{xcolor}
\usepackage{enumitem}
\usepackage{pifont}
\usepackage{graphicx}
\graphicspath{{../img}}
\usepackage{pgfplots}
\pgfplotsset{compat=1.18}
\usepackage{tabularx}
\usepackage{polynom}
\usepackage{setspace}
\usepackage{pdfpages}
\usepackage{tikz}
\usetikzlibrary{angles,quotes}


% ---------- General formatting ----------
\setlength{\parindent}{0pt}
\setlength{\parskip}{3pt}
\setlist[enumerate]{
    leftmargin=1.8em,
    itemsep=3pt,
    topsep=2pt,
    parsep=1pt
}

% ---------- Ruled writing lines ----------
\newcommand{\writingline}{%
    \noindent
    \makebox[\linewidth][l]{%
        \rule{0pt}{1.65em}%
        \leaders\hrule height 0.35pt\hfill
    }\par
}

\newcommand{\R}{\mathbb R}
\newcommand{\e}{{\rm e}}
\newcommand{\pobr}[1]{\left\langle#1\right\rangle}
\newcommand{\norm}[1]{\lVert #1 \rVert}
\newcommand{\abs}[1]{\lvert #1 \rvert}

\DeclareMathOperator{\xd}{d\!}
\DeclareMathOperator{\proj}{proj}

\title{}
\date{}

\newcommand{\template}{
\noindent
{
\textbf{Herramientas:} calculadora científica \hspace{\stretch{1}} Calificación \rule{4em}{.4pt} 
}

\Huge
\begin{enumerate}[leftmargin=*, itemsep=1.2em]
\item
 La forma completamente factorizada de la expresión \[x^6 - 63x^3 - 64\] es
 \begin{enumerate}[label = (\arabic*)]
\item
$(x-64)(x+1)$
\item
$(x^3 - 64)(x^3+1)$
\item
$(x+8)(x-8)(x+1)(x-1)$
\item
$(x-4)(x^2+4x+16)(x+1)(x^2-x+1)$
\end{enumerate} 
\end{enumerate}
\clearpage
}



\begin{document}

% ==========================================
% PAGE 1: Reading and reflection questions
% ==========================================

\begin{center}
    {\Large\bfseries Asumir la Responsabilidad de Tu Educación}

    \vspace{2pt}
    {\small\itshape Reflexión sobre el Primer Período de Calificaciones}
\end{center}

\textbf{¿Qué significa asumir la responsabilidad de tu educación?}

La educación es más que asistir a clases, completar tareas y
obtener calificaciones. Es una oportunidad para desarrollar conocimientos,
habilidades y formas de pensar que te servirán a lo largo de tu
vida. Tus profesores pueden guiarte, explicarte conceptos difíciles
y ofrecerte oportunidades para aprender. Sin embargo, no pueden hacer
el aprendizaje por ti. En última instancia, tu educación es algo en
lo que debes participar activamente y de lo cual debes hacerte responsable.

\textbf{Asumir la responsabilidad de tu educación significa tomar un papel
activo en tu propio aprendizaje.} Significa reconocer que tus elecciones,
hábitos y acciones importan. También significa ser honesto contigo mismo
sobre los desafíos que enfrentas, reconocer en qué puedes mejorar y
tomar medidas para alcanzar tus metas.

Asumir la responsabilidad no significa que todo esté bajo tu
control. Es posible que enfrentes circunstancias personales, responsabilidades
familiares, problemas de salud u otros desafíos que dificulten asistir
a clase o mantenerte enfocado. No necesitas culparte por circunstancias
fuera de tu control ni divulgar información privada para explicar tus
experiencias. En su lugar, concéntrate en comprender tu situación,
reconocer qué puedes influir y buscar ayuda cuando la necesites.

La responsabilidad también significa tomar la iniciativa cuando encuentras
dificultades. En matemáticas, cuando no entiendes un concepto, puedes
hacer preguntas, buscar aclaraciones o explicarle el concepto a otra
persona para poner a prueba tu comprensión. Cuando cometes un error,
puedes examinar tu razonamiento, identificar qué salió mal y aprender
de la experiencia en lugar de simplemente copiar la respuesta correcta.
Estas acciones demuestran que estás trabajando activamente para mejorar
en lugar de esperar a que alguien más resuelva cada dificultad por ti.

Asumir la responsabilidad también implica reflexionar honestamente sobre
tus hábitos y decisiones. Si faltas a clases, te cuesta participar o te
sientes desconectado, considera qué contribuyó a estas experiencias.
¿Qué factores estaban fuera de tu control? ¿Qué decisiones podrías haber
tomado de manera diferente? ¿Qué recursos o apoyo podrían ayudarte? El
propósito de la reflexión no es juzgarte ni buscar excusas, sino comprender
tus experiencias e identificar oportunidades de crecimiento.

El primer período de calificaciones es una oportunidad para reflexionar
sobre cómo han ido las cosas hasta ahora. No puedes cambiar el pasado,
pero puedes aprender de él. Al mirar hacia el próximo período de
calificaciones, piensa en una o dos acciones específicas y realistas que
puedas tomar para mejorar tu asistencia, participación o hábitos de
aprendizaje. En lugar de simplemente prometer esforzarte más, considera
qué harás realmente y qué apoyo podrías necesitar.

\textbf{Tu educación es un esfuerzo compartido, pero tu aprendizaje requiere
de tu participación. No tienes que hacerlo todo solo ni tienes que ser
perfecto. Lo que importa es que asumas un papel activo en tu propio progreso.}

\vspace{3pt}
\hrule
\vspace{4pt}

{\large\bfseries Reflexión Escrita: Mi Primer Período de Calificaciones}

Reflexiona sobre tus experiencias durante el primer período de calificaciones y
escribe una respuesta bien pensada que responda a las tres preguntas a continuación.
Utiliza ejemplos específicos según corresponda, pero no te sientas obligado
a divulgar información privada.

\begin{enumerate}
    \item \textbf{Asistencia:} ¿Cuáles fueron los motivos principales de tus
    ausencias durante el primer período de calificaciones? ¿Qué circunstancias
    o factores afectaron tu capacidad para asistir a clase con regularidad?

    \item \textbf{Participación en el aula:} ¿Cuáles fueron las razones principales
    por las que a veces te resultó difícil participar o mantenerte enfocado
    en clase? Considera los factores que afectaron tu atención, participación,
    motivación o disposición para pedir ayuda.

    \item \textbf{Responsabilidad y cómo seguir adelante:} ¿Qué significa para ti
    asumir la responsabilidad de tu educación? Mirando hacia atrás en el
    primer período de calificaciones, ¿cómo has asumido la responsabilidad de tu
    aprendizaje y en qué podrías mejorar? Identifica una o dos acciones específicas
    que tomarás durante el próximo período de calificaciones para mejorar tu
    asistencia, participación en clase o hábitos de aprendizaje. Si necesitas
    apoyo para realizar estos cambios, explica qué tipo de ayuda te serviría.
\end{enumerate}

Sé honesto y reflexivo. El propósito de esta tarea no es juzgar tu pasado
ni asignar culpas, sino comprender tus experiencias, reconocer tus
oportunidades de crecimiento y hacer un plan realista para salir adelante.
\clearpage

\noindent
{
\hspace{\stretch{1}} Calificación \rule{4em}{.4pt} 
}

\begin{center}
    {\Large\bfseries Mi Reflexión}
\end{center}

% Use any remaining space on page 1 for writing.
\vfill
\writingline
\writingline
\writingline

\writingline
\writingline
\writingline
\writingline
\writingline
\writingline
\writingline
\writingline
\writingline
\writingline
\writingline
\writingline
\writingline
\writingline
\writingline
\writingline
\writingline
\writingline
\writingline
\writingline
\writingline
\writingline
\writingline
\writingline
\writingline
\writingline
\writingline
\writingline
\end{document}